\section{Numerical Implementation}
\subsection{The General System}
The algorithm discussed in the previous work \citep{Chang2024} is repurposed here to handle various types of systems.
In the general nonlinear context, the governing system is
\begin{gather}\label{eq:eqv_sys}
    \left\{
    \begin{array}{l}
        \dot{\bbf}_j=-s_j\bbf_j+m_j\by_j, \\
        \br+\sum\bbf_j+\sum\bh_j=\mb{z}.
    \end{array}\right.
\end{gather}
Consider the external load may be augmented as in \eqsref{eq:a_final}, here we use $\mb{z}$ instead.
\subsubsection{Effective Stiffness}
The first expression is a first-order ODE (of $\bbf_j$) that can be integrated by the (implicit) trapezoidal rule, which is second-order accurate and A-stable,
\begin{gather}
    \bbf_{j,n+1}=\bbf_{j,n}+\dfrac{\Delta{}t}{2}\left(\dot{\bbf}_{j,n}+\dot{\bbf}_{j,n+1}\right).
\end{gather}
Expanding gives
\begin{gather}
    \bbf_{j,n+1}=\bbf_{j,n}+\dfrac{\Delta{}t}{2}\left(\left(-s_j\bbf_{j,n}+m_j\by_{j,n}\right)+\left(-s_j\bbf_{j,n+1}+m_j\by_{j,n+1}\right)\right),
\end{gather}
one may further rearrange to obtain
\begin{gather}\label{eq:discretised_c}
    \bbf_{j,n+1}=\dfrac{2-s_j\Delta{}t}{2+s_j\Delta{}t}\bbf_{j,n}+\dfrac{m_j\Delta{}t}{2+s_j\Delta{}t}\by_{j,n}+\dfrac{m_j\Delta{}t}{2+s_j\Delta{}t}\by_{j,n+1},
\end{gather}
in which subscripts $\left(\cdot\right)_n$ and $\left(\cdot\right)_{n+1}$ denote the corresponding quantity at $t_n$ and $t_{n+1}=t_n+\Delta{}t$.

Now assuming the equation of motion is satisfied at $t_{n+1}$, which is valid for a wide range of time integration methods (including the well-known Newmark method), the second expression is
\begin{gather}\label{eq:residual}
    \br_{n+1}+\sum\dfrac{2-s_j\Delta{}t}{2+s_j\Delta{}t}\bbf_{j,n}+\sum\dfrac{m_j\Delta{}t}{2+s_j\Delta{}t}\by_{j,n}+\sum\dfrac{m_j\Delta{}t}{2+s_j\Delta{}t}\by_{j,n+1}+\sum\bh_{j,n+1}=\mb{z}_{n+1}.
\end{gather}
Differentiation leads to the following revised effective stiffness $\bhat{K}_{n+1}$,
\begin{gather}\label{eq:revised_k}
    \bhat{K}_{n+1}=\bbar{K}_{n+1}+\sum\dfrac{m_j\Delta{}t}{2+s_j\Delta{}t}\ddfrac{\by_{j,n+1}}{\bu_{n+1}}+\sum\ddfrac{\bh_{j,n+1}}{\bu_{n+1}}.
\end{gather}
in which the conventional effective stiffness $\bbar{K}_{n+1}$ is
\begin{gather}\label{eq:effective_k}
    \bbar{K}_{n+1}=\ddfrac{\br_{n+1}}{\bu_{n+1}}=\mb{K}_{n+1}+\mb{C}_{n+1}\ddfrac{\bv_{n+1}}{\bu_{n+1}}+\mb{M}_{n+1}\ddfrac{\ba_{n+1}}{\bu_{n+1}}.
\end{gather}
In this work, we have considered two options: the trivial case $\bh_j=\mb{0}$ that corresponds to \eqsref{eq:v_final} and the non-trivial case $\bh_j=-\dfrac{m_j}{s_j}\by_j$ that corresponds to \eqsref{eq:u_final} and \eqsref{eq:a_final}.
For the latter,
\begin{gather}
    \sum\ddfrac{\bh_{j,n+1}}{\bu_{n+1}}=-\sum\dfrac{m_j}{s_j}\ddfrac{\by_{j,n+1}}{\bu_{n+1}}
\end{gather}
can be combined with the second term in \eqsref{eq:revised_k}.

It thus can be concluded from \eqsref{eq:revised_k} that to account for this energy dissipating action, only an additional matrix $\sum\dfrac{m_j\Delta{}t}{2+s_j\Delta{}t}\ddfrac{\by_{j,n+1}}{\bu_{n+1}}$, or $\sum\left(\dfrac{m_j\Delta{}t}{2+s_j\Delta{}t}-\dfrac{m_j}{s_j}\right)\ddfrac{\by_{j,n+1}}{\bu_{n+1}}$ if $\bh_j$ is present, needs to be added to the original effective stiffness $\bbar{K}_{n+1}$.
In certain cases, this additional matrix can be very efficiently formulated.
For example, if $\by_j\in\{\br_{\bu},\br_{\ba}\}$, then $\ddfrac{\by_{j,n+1}}{\bu_{n+1}}$ is simply the (possibly scaled) stiffness matrix $\mb{K}_{n+1}$ (or the mass matrix $\mb{M}_{n+1}$) that is already available in the original formulation.
In some other special cases, this matrix can be diagonal, for example, with $\by_j\in\{\bu,\bv,\ba\}$, then $\ddfrac{\by_{j,n+1}}{\bu_{n+1}}=\alpha\mb{I}$ is a diagonal matrix with possibly a non-unit scalar $\alpha$.
Adding such a diagonal matrix to $\bhat{K}_{n+1}$ is almost computationally trivial.
\subsubsection{Effective Load}
In \eqsref{eq:u_final} and \eqsref{eq:v_final}, simply
\begin{gather}\label{eq:eff_load_simple}
    \bz_{n+1}=\bp_{n+1}.
\end{gather}
In \eqsref{eq:a_final}, applying the same trapezoidal rule, one obtains
\begin{gather}
    \bq_{j,n+1}=\dfrac{2-s_j\Delta{}t}{2+s_j\Delta{}t}\bq_{j,n}+\dfrac{m_j\Delta{}t}{2+s_j\Delta{}t}\bp_n+\dfrac{m_j\Delta{}t}{2+s_j\Delta{}t}\bp_{n+1},
\end{gather}
thus
\begin{gather}\label{eq:eff_load}
    \begin{split}
        \bz_{n+1} & =\left(1-\sum\dfrac{m_j}{s_j}\right)\bp_{n+1}+\sum\bq_{j,n+1}                                                                                                                             \\
                  & =\sum\dfrac{2-s_j\Delta{}t}{2+s_j\Delta{}t}\bq_{j,n}+\sum\dfrac{m_j\Delta{}t}{2+s_j\Delta{}t}\bp_n+\left(1-\sum\dfrac{m_j}{s_j}+\sum\dfrac{m_j\Delta{}t}{2+s_j\Delta{}t}\right)\bp_{n+1}.
    \end{split}
\end{gather}
\subsubsection{A Reference Algorithm}
\algoref{algo:general_algo} presents an algorithm that integrates the proposed damping formulation into an existing iterative analysis workflow.
The subscript $\left(\cdot\right)_{n+1}$ is dropped for brevity.
Once convergence is achieved, $\bbf_{j,n}\leftarrow\bbf_j$ and $\bq_{j,n}\leftarrow\bq_j$ shall be committed.
When $\bh_j$ is an algebraic function of $\by_j$, there is usually no need to store it separately.
Even though $\bbf_j$ can be complex due to complex $s_j$ and $m_j$, as long as conjugate pairs of $s_j$ and $m_j$ are used, $\bbf=\sum\bbf_j$ and $\bq=\sum\bq_j$ are all real--valued.
The same also applies to scalars $\sum\dfrac{m_j\Delta{}t}{2+s_j\Delta{}t}$ and $\sum\dfrac{m_j}{s_j}$, and thus $\bhat{K}$ is real--valued.
\begin{breakablealgorithm}
    \setstretch{2}
    \caption{iteration body of solving dynamic systems with general damping}\label{algo:general_algo}
    \begin{algorithmic}
        \State \textbf{*Schema: Newmark + Trapezoidal}
        \State \textbf{Input}: $\bbar{K}$, $\bu$, $\bv$, $\ba$, $\br$, $\bp$ (quantities obtained in the conventional manner as if there is no general damping); damping history/parameters $\bbf_{j,n}$, $\by_{j,n}$, $\bq_{j,n}$, $m_j$, $s_j$
        \State \textbf{Output}: $\bu$
        \State \faMicrochip~compute energy dissipating force $\bbf_{j}=\dfrac{2-s_j\Delta{}t}{2+s_j\Delta{}t}\bbf_{j,n}+\dfrac{m_j\Delta{}t}{2+s_j\Delta{}t}\by_{j,n}+\dfrac{m_j\Delta{}t}{2+s_j\Delta{}t}\by_{j}$\Comment{\eqsref{eq:discretised_c}}
        \State \faMicrochip~compute effective effective stiffness $\bhat{K}=\bbar{K}+\sum\dfrac{m_j\Delta{}t}{2+s_j\Delta{}t}\ddfrac{\by_{j}}{\bu}+\sum\ddfrac{\bh_j}{\bu}$\Comment{\eqsref{eq:revised_k}}
        \State \faMicrochip~compute effective resistance $\bhat{r}=\br+\sum\bbf_j+\sum\bh_j$
        \State \faMicrochip~compute effective load $\bz$\Comment{\eqsref{eq:eff_load_simple} or \eqsref{eq:eff_load}}
        \State $\delta\bu=\bhat{K}^{-1}\left(\mb{z}-\bhat{r}\right)$
        \State update and return $\bu\leftarrow\bu+\delta\bu$
    \end{algorithmic}
\end{breakablealgorithm}
\subsection{Remarks}
\tabref{tab:special} summarises a few special cases.
\begin{table}[ht!]
    \centering\footnotesize\renewcommand{\arraystretch}{1.4}
    \caption{some noticeable formulations}\label{tab:special}
    \begin{tabular}{lccccll}
        \toprule
                                                  & $s_j$        & $m_j$        & $\by_j$                                       & $\bh_j$             & cost                                                                    & remarks                                                     \\ \midrule
                                                  & $\mathbb{C}$ & $\mathbb{C}$ & $\bu$                                         & $-m_j/s_j\bu$       & \faStar\faStarHalf*\faStar[regular]\faStar[regular]\faStar[regular]     & general damping                                             \\
                                                  & $\mathbb{C}$ & $\mathbb{C}$ & $\bv$                                         & $\mb{0}$            & \faStar\faStarHalf*\faStar[regular]\faStar[regular]\faStar[regular]     & identical to \citep{Chang2024}                              \\
                                                  & $\mathbb{C}$ & $\mathbb{C}$ & $\ba$                                         & $-m_j/s_j\ba$       & \faStar\faStarHalf*\faStar[regular]\faStar[regular]\faStar[regular]     & general damping                                             \\ \midrule
                                                  & $\mathbb{R}$ & $\mathbb{R}$ & $\dot{\br}_{\bu}$ (linear: $\mb{K}\bv$)       & $\mb{0}$            & \faStar\faStar\faStar\faStarHalf*\faStar[regular]                       & identical to \citep{Huang2019}                              \\
        \eqsref{eq:v_final}, \secref{sec:v_dep_k} & $\mathbb{C}$ & $\mathbb{C}$ & $\dot{\br}_{\bu}$ (linear: $\mb{K}\bv$)       & $\mb{0}$            & \faStar\faStar\faStar\faStar\faStar[regular]                            & complex extension of \citep{Huang2019}                      \\
        \secref{sec:v_dep_m}                      & $\mathbb{C}$ & $\mathbb{C}$ & $\dot{\br}_{\ba}$ (linear: $\mb{M}\dot{\ba}$) & $\mb{0}$            & \faStar\faStar\faStar\faStar\faStar                                     & equiv. to \secref{sec:v_dep_k} if free vibration            \\ \midrule
        \eqsref{eq:u_final}                       & $\mathbb{C}$ & $\mathbb{C}$ & $\br_{\bu}$                                   & $-m_j/s_j\br_{\bu}$ & \faStar\faStar\faStarHalf*\faStar[regular]\faStar[regular]              & alt. form, equiv. to \secref{sec:v_dep_k}                   \\
        \eqsref{eq:a_final}                       & $\mathbb{C}$ & $\mathbb{C}$ & $\br_{\ba}$                                   & $-m_j/s_j\br_{\ba}$ & \faStar\faStar[regular]\faStar[regular]\faStar[regular]\faStar[regular] & alt. form, equiv. to \secref{sec:v_dep_k}                   \\ \midrule
        \multicolumn{7}{p{15cm}}{The reported implementation cost is intended only to indicate the relative ranking among the formulations in the table, and does not represent a quantitative or absolute measure of computational expense. More stars indicate higher implementation cost.} \\ \bottomrule
    \end{tabular}
\end{table}
The proposed damping formulations can be applied at both the global model and element levels.
Consequently, the energy dissipating forces $\bbf$ and $\bh$ can be flexibly formulated using either node-based or element-based representations.

One shall notice that the proposed formulations have the following advantages.
\begin{enumerate}
    \item It retains the size of the original problem/system.
    \item The parameters $s_j$ and $m_j$ accept complex values, providing extra flexibility customising the desired damping response.
    \item State determination of additional damping history $\bbf_j$ is constant and only involves matrix--vector multiplications.
    \item The consistent tangent stiffness that allows quadratic convergence rate can be trivially assembled, for example, using different scalar weights in weighted summation if $\by_j\in\left\{\br_{\bu},\br_{\bv},\br_{\ba}\right\}$ or assembling additional diagonal matrices if $\by_j\in\left\{\bu,\bv,\ba\right\}$.
          In problems where the mass matrix $\mb{M}$ is constant and/or lumped, the extra computational efforts required would be minimized.
    \item The reciprocal formulations are equivalent mathematically.
          In the absence of other energy dissipating actions (such as conventional viscous damping), one form can be converted into another.
          Instead of scaling $\mb{K}$, it is possible to scale $\mb{M}$ in the equivalent form.
          In global time integration methods that express the equation of motion in terms of $\bu$, $\mb{M}$ is already scaled by a scalar factor; therefore, the proposed contribution can be included with one additional scalar multiplication.
\end{enumerate}

The preceding analyses assume no external damping actions.
Consequently, the derived loss stiffnesses and loss ratios hold strictly in the absence of additional velocity dependent energy dissipation mechanisms, such as conventional viscous damping.
Nonetheless, a non-trivial $\mb{C}$ is maintained in \eqsref{eq:effective_k} to preserve structural generality.